\section{Inverting the scalar pair kernel}
\label{sec:inverse}

The cavity equations couple errors carried by different pairs of
leaves. To solve them, we need an inverse whose bound survives refinement
of the hierarchy and the small masses used in the comparison. The
proposition below provides that bound for arbitrary fixed old trees;
its proof keeps the signed cancellations before taking absolute values.

A pair function assigns a deterministic number to each pair and each
allowed extension of the fixed old-label tree. In the applications,
that number is a positive-law expectation of an overlap test. The
operator sums below vary the extended topology using signed allocation
coefficients; the scalar expectations in its kernel use the positive
law on each such topology.

For a finite resolution with terminal mass one, choose mass cuts
\(0<p_0<\cdots<p_{R+1}=1\). A \emph{sheet} \(r\) consists of
pairs in the same block at cut \(p_r\) and different child blocks
at cut \(p_{r+1}\); write \(K_r\) for their common scalar time.
The resolution records these relations to all fixed old labels, so a
pair function need not be constant on a sheet. The ordinary absolute
attachment rule means the signed rule of Section~\ref{sec:foundations}
with each coefficient replaced by its absolute value. In particular,
it retains both the density for a new fork and the atom at an existing
fork, as well as terminal coincidences.

\begin{proposition}[Pair inverse with absolute attachment domination]
\label{inv:pair-inverse}
The kernel \eqref{eq:A8} is considered on extensions of any \emph{fixed} finite list of old labels. A pair function \(u_{ab}\) may depend on the topology with those labels but ignores any unused extra labels (is pulled back by restriction). Indices in this section used as pair arguments are distinct leaves. Normalize terminal mass to 1; other bounded terminal masses reduce to this by mass rescaling, scaling cumulative covariances by the square of the mass. Take a monotone step scalar path \(K\), nonnegative bounded-variation step multiplier \(b\), and retain pair levels in a terminal interval of the hierarchy (the whole interval is allowed). Suppose at these levels
\[
 1-b_r\mathbf E\chi_r^2\ge e>0,\qquad \chi_r=V_{zz}(K_r,Z_{K_r}).
\]
Bounds called polynomial in this section are in \(1+K_{\rm top}+\|b\|_\infty+\operatorname{var}b+e^{-1}\), independent of the number of steps. On the retained levels
\[
 u_{ab}\ \longmapsto\ u_{ab}-\sum_{cd\ {\rm retained}}{\bf M}_{ab,cd} b_{cd}u_{cd}
                                                                     \tag{H1}\label{eq:H1}
\]
has an inverse on each finite step resolution. Each output is an integral of inputs over allocations of up to two more labels (some arguments can also be reused), with weights dominated by a polynomial times the ordinary \emph{absolute} unrestricted attachment rules to the old tree including \(ab\). Restrictions and bounded weights can depend on the topology. This formulation in particular controls new split densities, allowing coincidences and ties to existing forks. It applies with \(ab\) using old labels too. We only need symmetric arrays (on antisymmetric ones the subtracted sum is zero).

\end{proposition}

The proof separates the scalar calculation from the construction of the
inverse. We first resolve pair arrays into two sectors and state the
scalar response identities they require. With those identities in hand,
we construct the inverse and prove its attachment bound. The detailed
scalar verification follows afterward.

\subsection{Pair resolutions and scalar response}

Initially keep the incoming root power \(p_0>0\) and the clocks
strictly ordered. We will remove both restrictions after obtaining
bounds independent of the resolution.

The scalar time profile is \(\alpha=p_i\) on \((K_{i-1},K_i)\), where
\(K_{-1}=0\). Add a mass-one terminal interval ending at
\(K_{R+1}\); it leaves all earlier field derivatives unchanged.
For a function of a leaf, \(P_k\) replaces that leaf by a signed
allocation within its \(p_k\)-block and divides the result by \(p_k\).
The signed row-sum rule therefore gives \(P_k1=1\), and
\(P_{R+1}=I\). Put
\(E_0=P_0\) and \(E_k=P_k-P_{k-1}\) for \(k\ge1\).
These differences separate successive block resolutions. When a pair
lies on sheet \(t\), resolve each endpoint only within its own child:
the coarsest operator is then
\(\widetilde E^t_{t+1}=P_{t+1}\), followed by
\(\widetilde E^t_i=E_i\) for \(i>t+1\).

\paragraph{The additive and row-zero sectors}
For the counting calculation alone, treat the symbols \(p_i\) as integer
block sizes with sufficiently large integer ratios \(p_i/p_{i+1}\).
Thus this temporary integer tree has decreasing block sizes; the
increasing real masses are substituted only after the counting
identities have been proved. On the integer tree use symmetric
off-diagonal pair arrays, with inner product summed over unordered
pairs.

Nested block averages on this tree satisfy
\[
 P_iP_j=P_{\min(i,j)},\qquad
 E_kE_l=\mathbf1_{\{k=l\}}E_k,\qquad \sum_{k=0}^{R+1}E_k=I.
\]
The coefficient argument below transfers these identities to real masses.

Fix a leaf function \(\epsilon\) satisfying \(E_k\epsilon=\epsilon\).
A \emph{sheet-additive} array has value
\(f_t(\epsilon_a+\epsilon_b)\) when \(ab\) lies on sheet \(t\).
The coefficients \(f_t\) may vary with the sheet. Summing over the
second leaf in a row on sheet \(t\) gives
\(d_t^k f_t\epsilon_a\), where
\[
 d_t^k=p_t-p_{t+1}+p_t{\bf1}_{k\le t}-p_{t+1}{\bf1}_{k\le t+1},\qquad \mu_t^k=-d_t^k .
\]
On this integer tree, its squared norm divided by
\(\sum_a\epsilon_a^2\) is \(\sum_t d_t^k f_t^2\).
The kernel \({\bf M}\) preserves the arrays of this form and acts on
the level coefficients \(f_t\) by a matrix independent of
\(\epsilon\). To see this, fix the output pair \(ab\). Tree symmetry
makes the sum of a one-leaf input constant on each block away from the
paths to \(a\) and \(b\). The resulting leaf operators are block
projections along those two paths. On the range of \(E_k\), each such
projection is either zero or the identity, so the answer is a multiple
of \(\epsilon_a+\epsilon_b\).

For a general pair array, write
\((R_tu)_a=\sum_{b:\,ab\ {\rm on\ sheet}\ t}u_{ab}\).
The additive projection is explicitly
\[
 (\Pi_{\rm add}u)_{ab}\big|_t
 =\sum_k\frac{(E_kR_tu)_a+(E_kR_tu)_b}{d_t^k}.
\]
Subtracting this projection leaves an array whose row sum on each
sheet is zero. We call this the \emph{row-zero} part. For a
row based at a leaf in a \(p_t\)-block, the sum is over all other
children of that block together; it need not vanish separately in each
child.

\paragraph{Real masses and arbitrary old-label topologies}
The transfer to real masses concerns coefficients of moment expressions,
not Gaussian expectations as functions of an integer replica count.
Introduce a formal variable for each pair or four-leaf moment, recording
all leaf coincidences and overlap levels. Include products of pair
moments as separate expressions. Also leave every input value as a
formal variable indexed by its complete genealogy with the old labels.
After expanding a block projection or kernel action, the coefficient of
each such expression is a rational function of the \(p_i\).

The counts in these coefficients are falling powers: when adding a leaf,
one either uses an occupied child or chooses a new child. These are
exactly the coefficients in the real-power rule of PE, Section~2.
To verify a coefficient identity, clear denominators and regard the
branching ratios as independent variables. Integer counting shows that
the resulting polynomial vanishes for every sufficiently large integer
choice of these variables, also if even ratios are required. Applying
the one-variable polynomial identity successively in each variable
shows that the polynomial is identically zero. We may therefore
substitute the increasing real masses.

In particular, the additive projection is obtained by taking a row sum
on sheet \(t\), applying \(E_k\), dividing by \(d_t^k\), and adding
the result at the two endpoints. The denominators are nonzero for
strictly increasing real masses. Subtracting these projections gives
the stated row-zero remainder. Because the input variables retain
their entire old-label genealogy, these identities apply to arbitrary
pair functions in the proposition, not just sheet-constant functions.
The signed real-mass calculation uses these coefficient identities;
it makes no assertion about a vector space of negative dimension.

The decomposition reduces inversion to two scalar calculations. On the
row-zero sector, the kernel is diagonal in the two child resolutions.
On the additive sector, it is a diagonal operator minus a positive
Hessian kernel integrated against sheet masses. The next lemma gives
both actions and the estimates that will cancel block denominators.

\begin{lemma}[Scalar response on the two sectors]
\label{inv:sector-response}
Use the finite resolution above, and put
\(m_t=V_z(K_t,Z_{K_t})\), \(\chi_t=V_{zz}(K_t,Z_{K_t})\).
Let \(\mathbf E_t\) condition on the scalar prefix through \(K_t\).
On the row-zero sector, different sheets do not mix. On sheet \(r\),
the multiplier in the child resolutions
\(\widetilde E_i^r\otimes\widetilde E_l^r\), \(i,l\ge r+1\), is
\[
 g_r^{il}=\mathbf E[(\mathbf E_r\chi_{i-1})
                         (\mathbf E_r\chi_{l-1})].
\]
On the additive sector associated with \(E_k\), the action on
sheet coefficients is
\[
 ({\bf M}^k f)_r=g_r^k f_r-\sum_t B^k_{rt}\mu_t^k f_t,\qquad
 g_t^k=\mathbf E[\chi_t\chi_{\max(k-1,t)}].
 \tag{H2}\label{eq:H2}
\]
Here \(A=K_{k-1}\), with \(A=0\) for \(k=0\), and
\[
 B^k_{rt}=\mathbf E\left[
   \operatorname{Hess}_{z,\alpha}V(A,Z_A)[\beta_r,\beta_t]\right],
 \qquad
 \beta_t=\begin{cases}
 (m_t,0),&t<k,\\
 (0,\delta_{K_t}),&t\ge k.
 \end{cases}
\]
The Hessian varies the starting field and the subsequent mass profile.
The prefix values \(m_t\), \(t<k\), are held fixed during this
differentiation, and \(\delta_{K_t}\) denotes an impulse in the
subsequent profile. Expectation is taken after evaluating the Hessian.
The matrix \(B^k\) is positive semidefinite. Restricting either
action to retained levels discards the other input and output sheets.

The kernels \(B^k,g^k\) are bounded and Lipschitz in their evaluation
times with polynomial constants. Their changes with the resolution
index satisfy
\[
 \|B^{k+1}-B^k\|_\infty+\|g^{k+1}-g^k\|_\infty
 \lesssim_{\rm pol}p_k(K_k-K_{k-1}).
\]
A successive difference of \(g_r^{il}\) in index \(i\) is bounded
by \(Cp_i(K_i-K_{i-1})\); a mixed difference is bounded by the
product of the corresponding two bounds. Moreover
\(0\le g_r^k,g_r^{il}\le\mathbf E\chi_r^2\).
These formulas apply with arbitrary old-label topologies under the
coefficient convention just described. For \(k=0\), their Hessian
kernel also satisfies
\[
 |B^0_{rt}|\lesssim C(1+K_{R+1})^C K_{\min(r,t)}
\]
when the earlier evaluation time is small.
\end{lemma}

We prove the lemma after constructing the inverse. Its role is now
explicit: positivity controls the additive resolvent, while its weighted
differences cancel the normalizations of the block projections.

\subsection{Construction and attachment bound for the inverse}

\begin{proof}[Proof of Proposition~\ref{inv:pair-inverse}]
If there are no retained sheets, the assertion is immediate. Otherwise
use the decomposition above. By Lemma~\ref{inv:sector-response}, the
gap hypothesis implies
\[
 \ell_t^k=1-b_tg_t^k\ge e,\qquad
 \ell_t^{il}=1-b_tg_t^{il}\ge e
\]
on each retained sheet. We invert the two sectors, combine their block
projections, and only then estimate their attachment coefficients.

\paragraph{The additive resolvent}
Fix \(k\) and retain only the permitted rows and columns. Abbreviate
\(B=B^k\) and
\(w=\operatorname{diag}(\mu_t^k b_t/\ell_t^k)\).
The diagonal matrix \(w\) is nonnegative. Consequently
\(C=(I+B w)^{-1}B\) is symmetric and satisfies
\(0\preceq C\preceq B\), as is seen by writing it as
\(B^{1/2}(I+B^{1/2}wB^{1/2})^{-1}B^{1/2}\).
The additive-sector inverse is
\((\ell^k)^{-1}I+A^k\operatorname{diag}(\mu^k)\), where
\(A^k_{rt}=-C_{rt}b_t/(\ell_r^k\ell_t^k)\).
The bound \(C\preceq B\) controls the diagonal entries of \(C\), and
positivity gives \(|C_{rt}|\le\sqrt{C_{rr}C_{tt}}\). Thus all entries
are polynomially bounded. The identity \(C=B-BwC\), together with
\(\sum_t w_{tt}=O_{\rm pol}(1)\), transfers the evaluation-time
Lipschitz bounds for \(B\) to \(C\) and then to \(A\). For \(A\),
the relevant distance is \(|\Delta K|+|\Delta b|\).
Successive resolution indices satisfy the stronger estimate
\[
 \|A^{k+1}-A^k\|_\infty\lesssim_{\rm pol}
 p_k(K_k-K_{k-1}+|b_k-b_{k-1}|+{\bf1}_{k=\min(\mathrm{retained})})
\]
where \(b_{-1}=b_0\). To prove this, note that changing \(k\) to
\(k+1\) changes the sheet measure \(\mu\) by
\(p_k(\delta_k-\delta_{k-1})\), with unretained atoms omitted.
Use primes for the matrices at \(k+1\). Their resolvent identity is
\[
 C'-C=(I-C'w')(B'-B)(I-wC)-C'(w'-w)C .
\]
If both atoms are retained, their contributions to
\(C'(w'-w)C\) are subtracted before taking absolute values. The
evaluation Lipschitz estimates then supply
\(p_k(K_k-K_{k-1}+|b_k-b_{k-1}|)\). Changes in the kernel and in the
diagonal gaps use the preceding bound with the same factor \(p_k\).

At the first retained sheet, the lower atom is absent. Its partner has
mass exactly \(p_k\), which accounts for the indicator in the bound.
There is only one such boundary term. It still has the factor needed
to cancel the block-average denominator.

\paragraph{Absolute attachment domination}
Fix an output pair \(ab\) on a retained sheet \(r\). First apply the
diagonal multiplier \(1/\ell_r^{il}\) in the two child resolutions.
Then add the row correction
\[
 -\sum_{k,t}A^k_{rt}\{(E_k\,{\rm row}_t\,u)_a+(E_k\,{\rm row}_t\,u)_b\},
 \qquad ({\rm row}_t\,u)_c=\sum_{d:\,cd\ {\rm on\ sheet}\ t}u_{cd}.
\]
This is the inverse on both parts of the decomposition. On a row-zero
array, the correction vanishes. In each one-leaf summand of an additive
array, the other endpoint is constant within its child and therefore
has child index \(r+1\). The child multiplier on the summand is thus
\(1/\ell_r^k\). The row sum contributes
\(d_t^k=-\mu_t^k\), and the minus sign in the correction gives the
remaining \(A^k\operatorname{diag}(\mu^k)\) part of the inverse.

We now combine the block projections before estimating them. Since
\(E_0=P_0\), \(E_k=P_k-P_{k-1}\), and \(P_{R+1}=I\), the exact
one-variable Abel identity is
\[
 \sum_{k=0}^{R+1}A^k_{rt}E_k
  =A^{R+1}_{rt}I+\sum_{k=0}^{R}(A^k_{rt}-A^{k+1}_{rt})P_k .
\]
Thus the only terminal term is an identity operator. In every other
term, the coefficient difference has a factor \(p_k\), canceling the
denominator of \(P_k\).

Here is the corresponding pointwise comparison with attachment
densities. Write \(c\sim_k a\) if \(c\) belongs to the \(p_k\)-block
of \(a\), and let \(\mathcal S(dc)\) denote the ordinary signed
attachment of \(c\) to the old tree, including \(ab\). By definition,
\[
 (P_kF)_a=\frac1{p_k}\int {\bf1}_{\{c\sim_k a\}}F_c\,\mathcal S(dc).
\]
A subsequent row sum attaches \(d\) to the tree that now includes
\(c\); denote that signed rule by \(\mathcal S_c(dd)\).
Relative to these two successive attachment rules, the nonterminal
correction at the endpoint \(a\) has coefficient
\[
 \kappa_a(c,d)=
 -\sum_{k=0}^{R}\sum_{t\ {\rm retained}}
 \frac{A^k_{rt}-A^{k+1}_{rt}}{p_k}
 {\bf1}_{\{c\sim_k a\}}
 {\bf1}_{\{cd\ {\rm on\ sheet}\ t\}} .
\]
For a fixed allocation, the pair \(cd\) belongs to at most one sheet.
Put
\(\rho_k=K_k-K_{k-1}+|b_k-b_{k-1}|
 +{\bf1}_{k=\min(\mathrm{retained})}\).
The weighted difference estimate therefore gives
\[
 |\kappa_a(c,d)|\lesssim_{\rm pol}
 \sum_{k=0}^{R}\rho_k{\bf1}_{\{c\sim_k a\}}
 \le K_R+\operatorname{var}b+1 .
\]
This bounds the coefficient at each placement relative to the
unrestricted absolute rule
\(|\mathcal S|(dc)|\mathcal S_c|(dd)\).
The endpoint \(b\) has the same bound. The terminal identity term
requires just the row attachment of \(d\), with a polynomially bounded
coefficient.

For the child multiplier, abbreviate \(D_{il}=1/\ell_r^{il}\),
\(i,l=r+1,\ldots,R+1\). Set
\(\Delta_iD_{il}=D_{il}-D_{i+1,l}\) and
\(\Delta_lD_{il}=D_{il}-D_{i,l+1}\). Applying the same Abel identity
in both positions gives all endpoint terms explicitly:
\[
\begin{split}
 \sum_{i,l=r+1}^{R+1}D_{il}\,
       \widetilde E_i^r\otimes\widetilde E_l^r
 ={}&D_{R+1,R+1}I\otimes I\\
 &+\sum_{i=r+1}^{R}\Delta_iD_{i,R+1}P_i\otimes I
  +\sum_{l=r+1}^{R}\Delta_lD_{R+1,l}I\otimes P_l\\
 &+\sum_{i,l=r+1}^{R}\Delta_i\Delta_lD_{il}P_i\otimes P_l .
\end{split}
\]
There is no separate lowest-child average: it is included in the
difference sums. With \(\Delta K_i=K_i-K_{i-1}\), the estimates for
\(g_r^{il}\) and the gap \(\ell_r^{il}\ge e\) give
\[
\begin{split}
 |\Delta_iD_{i,R+1}|&\lesssim_{\rm pol}p_i\Delta K_i,\qquad
 |\Delta_lD_{R+1,l}|\lesssim_{\rm pol}p_l\Delta K_l,\\
 |\Delta_i\Delta_lD_{il}|
       &\lesssim_{\rm pol}p_i p_l\Delta K_i\Delta K_l .
\end{split}
\]
The single differences cancel one block denominator. The coefficient
of the two-attachment term is bounded pointwise by
\[
 \sum_{i,l=r+1}^{R}
 \frac{|\Delta_i\Delta_lD_{il}|}{p_i p_l}
 {\bf1}_{\{c\sim_i a\}}{\bf1}_{\{d\sim_l b\}}
 \lesssim_{\rm pol}
 \left(\sum_{i=r+1}^{R}\Delta K_i\right)^2 .
\]
The identity term uses no new label; the two edge sums use one; and the
last sum uses at most two. All are thus bounded by the claimed
unrestricted absolute attachment rules.

The allocation order identities in Section~\ref{sec:foundations} allow
the block and old-topology restrictions to be retained throughout these
equalities. We discard them only in the pointwise inequalities above.
The summable coefficients depend on total clock length, total variation
of \(b\), and the single retained boundary, not on the number of steps.
Existing-fork atoms and coincidences keep their ordinary attachment
weights. This also explains why the argument controls new split
densities, rather than only the total mass of the inverse.

\paragraph{Limits of finite resolutions}\label{par:inverse-limits}
At a fixed finite resolution, let strictly separated clocks approach
flat clocks and then let \(p_0\downarrow0\). The combined coefficients
above stay bounded after their block denominators have been canceled.
Taking subsequences of these finitely many coefficients preserves the
inverse identity and its absolute bounds. A strict gap persists after
an arbitrarily small weakening.

The same resolution operators can be applied to continuum placements:
retain the indicators of membership in the blocks at each resolution
cut, while keeping the exact old-label positions. To check this, insert
empty steps between the cuts. Integer counting of additional slots
depends only on their coarse genealogy with the old labels, so refining
a group of children does not change that count. The real-mass
coefficient identity then gives the same conclusion. Consequently the
pointwise comparison with the unrestricted attachment measure remains
valid for continuum placements, including atoms at old forks.

For a nonstep path, apply the inverse at finite resolutions and fixed
system size. PE, Proposition~2.14 gives convergence of positive
expectations at each fixed placement, preserving prescribed cuts and
excluding only a null set of new placements. The inverse kernels need
not themselves converge: their common absolute majorant bounds the
error made by replacing a bounded step-path observable with its limit,
and dominated convergence makes that error tend to zero. This is the
form used in the applications.

Thresholds, prescribed levels, and specified sides of jumps can be
included in every approximation. Sample simultaneous path values and,
if needed, weaken the strict gap slightly; at fixed bounded clocks,
the scalar PDE derivatives and moments are uniformly continuous under
\(L^1\) approximation of the mass profile in common time. Finally,
for sheet-only inputs the \(k=0\) formula \eqref{eq:H2} passes to the
limit with \(\mu^0=2\,dv\).
\end{proof}

\subsection{Verification of the scalar response formulas}

\begin{proof}[Proof of Lemma~\ref{inv:sector-response}]
We first calculate the row-zero action by conditional independence.
The additive action then comes from a colored covariance variation.
Finally, explicit scalar Hessian kernels give positivity and the
weighted estimates used in the inverse.

\paragraph{The row-zero action}
Fix an input \(p_t\)-block on the integer tree and condition on
its shared scalar prefix. Within a child define
\(C_{ac}=\mathbf E_t[\tau_a\tau_c]-m_t^2\), and set \(C_{ac}=0\)
when the two leaves are in different children. If the output leaves
also lie in different children of the input block, conditional
independence gives the following four-spin moment for an input pair
\(cd\) on sheet \(t\):
\[
 m_t^4+m_t^2(C_{ac}+C_{ad}+C_{bc}+C_{bd})+C_{ac}C_{bd}+C_{ad}C_{bc}.
\]
The constant and one-end terms vanish when summed against a row-zero
input. The disconnected subtraction in \eqref{eq:A8} vanishes for
the same reason: the input pair moment is constant on this block and
sheet. The two remaining products give \(C\otimes C\); the two
orderings of the input pair cancel the factor \(1/2\) in
\eqref{eq:A8}. If the output does not occupy two different children of
this same block, each kernel term instead depends on at most one input
end and its sum is zero. This proves that different sheets do not mix.

The conditional matrix \(C\) is diagonal in the child resolutions. Its
values at successive overlap levels are
\(L_j=\mathbf E_t m_j^2-m_t^2\), with \(L_t=0\) and
\(L_{R+1}=1-m_t^2\). The eigenvalue at index \(i\) is
\(\sum_{j=i}^{R+1}p_j(L_j-L_{j-1})=\mathbf E_t\chi_{i-1}\),
by the conditional source-Hessian formula. Multiplying the two
conditional eigenvalues and averaging the prefix gives \(g_t^{il}\).
Finally, \(C\) has the same constant row sum in each child. Its tensor
action and the child projections therefore preserve the row-zero
condition used above.

This calculation transfers to real masses by the coefficient rule
above: the entries of \(C\) are differences of pair moments, and
products under a shared prefix are four-leaf moments with independent
descendant branches. Keeping the complete old-label genealogy in the
input variables preserves the formula for arbitrary pair functions.

\paragraph{A colored covariance calculation}
To compute the matrix in \eqref{eq:H2}, it suffices to compute its
quadratic form on an \(E_k\) sector. For \(k\ge1\), color half of the
children at depth \(k\) in each parent \(+\) and half \(-\), and give
every descendant the sign \(\epsilon\) of its child. This leaf function
lies in the range of \(E_k\). For real powers the corresponding
expression has two child integrals, each raised to
\(p_{k-1}/(2p_k)\). For \(k=0\), use the constant \(+\) function.

Multiply each Gaussian increment on a leaf's path by
\(1+\theta\epsilon h_j\) at depth \(j\), including increments
inherited from above the colored fork. Choose
\(h_j(K_j-K_{j-1})=f_j-f_{j-1}\), where \(f_{-1}=0\), and choose
any terminal value \(f_{R+1}\). Strictly separated clocks make this
choice possible. Write
\(G_j=\sum_{s\le j}h_s^2(K_s-K_{s-1})\), \(G_{-1}=0\), and
\(s_t=p_{t+1}-p_t\).

Differentiate the log root partition divided by \(p_0\) twice at zero. Covariance acceleration contributes
\[
 G_{R+1}-\sum_{t\ge k}s_t G_t\mathbf E m_t^2-p_k G_{k-1}\mathbf E m_{k-1}^2
\]
with the last term omitted for \(k=0\). The remaining curvature is
\(\sum_t d_t^k f_t({\bf M}^k f)_t\). Indeed, on sheet \(t\) the
first covariance variation is \((\epsilon_a+\epsilon_b)f_t\), and the
coefficient of \(\theta^2\) is \(\epsilon_a\epsilon_bG_t\). The
spin-square diagonal is constant and contributes no covariance Hessian.
On an integer complete tree, ordinary differentiation with
\(\sum_a\epsilon_a^2=p_0\) gives the claimed contraction, pattern by
pattern in the moments.

For real powers, carry the inherited \(+\) and \(-\) inputs as a
two-component Gaussian field until the color separation, where the two
child powers are multiplied. Gaussian covariance differentiation and
the labeled product rule assign two hits to the same or different
children with the same falling-power coefficients as in the integer
calculation. Only shared increments enter a covariance. At
\(\theta=0\), every normalized transition has the unperturbed scalar
law, so the moment factors are the ones required above.

The disconnected product comes from differentiating the root logarithm.
Its two pair sums are independent symbolic sums; on a common extended
topology, each moment is evaluated on its own restriction. The
allocation identity from Section~\ref{sec:foundations} preserves their
product. All coefficients therefore agree by the real-mass argument,
with singular Gaussian covariance matrices handled by continuity. This
identifies the derivative with the contraction before we evaluate it.

\paragraph{Evaluation of the colored Hessian}
We now use the increasing real masses and evaluate the derivative under
the unperturbed positive scalar law. Continue to write
\(A=K_{k-1}\), and temporarily put \(u=m_{k-1}\),
\(F=f_{k-1}\), \(p=p_k\), and \(\chi=\chi_{k-1}\).
When \(k=0\), take \(A=0\), \(F=0\), and evaluate \(u,\chi\) at
the initial state. Let \(\mathcal Z\) be the sum of the Gaussian inputs before
color separation, each weighted by \(h_j\), \(j<k\).

For a positive-color child, let \(V_\theta\) be the continuation whose
times, including its starting and ending times, have moved to
\(K_t+2\theta f_t+\theta^2G_t\). Dots below mean derivatives at
\(\theta=0\). For \(k>0\), the first derivatives from the two
opposite colors cancel pointwise before the fork. The second derivative
of the normalized root value is consequently
\[
 \mathbf E[\ddot V+2\mathcal Z\dot V_z+\mathcal Z^2\chi] .
\]
For \(k=0\) it is the root continuation derivative with \(\mathcal Z=0\).

Set \(\eta=2\sum_{t\ge k}s_tf_t\delta_{K_t}\), \(J=D_\alpha V_A[\eta]=\sum_{t\ge k}s_tf_t\mathbf E_A m_t^2\). Differentiating the moved times gives (at fixed starting field)
\[
\begin{split}
 \dot V&= f_{R+1}-F(\chi+p u^2)-J,\\
 \ddot V&=G_{R+1}-G_{k-1}(\chi+p u^2)-\sum_{t\ge k}s_t G_t\mathbf E_A m_t^2
         -2\sum_{t\ge k}s_t f_t^2\mathbf E_A\chi_t^2\\
      &\qquad+4F^2\partial_{AA} V-4F\partial_A J+
                       \operatorname{Hess}_{\alpha\alpha}V_A[\eta,\eta].
\end{split}
\]
To account for the terms, moving a step location changes the profile by
minus its jump times the displacement. Its first response is
\(\mathbf E_A m_t^2/2\), whose time derivative is
\(\mathbf E_A\chi_t^2/2\). These give \(-J\) in the first line and
the two sums over \(t\ge k\) in the second. The remaining terms
differentiate the starting time, the annealed terminal constant, and
the profile twice.

The impulses can also be obtained by integrating over the short
intervals swept out by the moving steps. The Hessian kernels below are
continuous when the impulse times meet. At fixed bounded clocks, the
continuation and its needed derivatives are continuous under
\(L^1\) profile perturbations, by their bounded-coefficient diffusion
equations. This justifies the impulse limit. Extend the initial
constant-mass interval slightly when differentiating the starting time;
strict clock separation keeps all other steps outside that interval.

The prefix is sampled under its normalized positive law, rather than
under the original Gaussian measure. Multiplying its successive density
factors gives a telescoping exponential. Gaussian integration by parts
therefore differentiates both the endpoint test and this density. For
\(Y\) depending on the endpoint field, the first two identities are
\[
\begin{split}
 \mathbf E[\mathcal ZY]&=\mathbf E(FY_z+HY),\\
 \mathbf E[(\mathcal Z^2-G_{k-1})Y]&=\mathbf E[F^2Y_{zz}+2FHY_z+(H^2+N')Y],\\
 H&=pF u-\sum_{t<k}s_t f_t m_t,\qquad
 N'=pF^2\chi-\sum_{t<k}s_t f_t^2\chi_t .
\end{split}
\]
Here \(H\) and \(N'\) are the first and second contributions from
differentiating the tilted density. The starting-time PDE gives
\(4V_{AA}=\chi_{zz}+2p\chi^2+4p u\chi_z+4p^2u^2\chi\) and
\(-4J_A=2J_{zz}+4p uJ_z\). Insert these and the two integration-by-parts
identities into the second derivative. After subtracting the covariance
acceleration displayed above, the contraction becomes
\[
 -p F^2\mathbf E\chi^2-\sum_{t<k}s_t f_t^2\mathbf E(\chi\chi_t)
 -2\sum_{t\ge k}s_t f_t^2\mathbf E\chi_t^2
 +\mathbf E\left[\operatorname{Hess}_{z,\alpha}V_A[\xi,\xi]\right],
\]
where \(\xi=(pF u+\sum_{t<k}s_t f_t m_t,\eta)\), and the Hessian is
evaluated as a bilinear form. For \(t<k\),
\(\mu_t^k=s_t+p_k{\bf1}_{t=k-1}\); for \(t\ge k\),
\(\mu_t^k=2s_t\). Thus \(\xi=\sum_t\mu_t^k f_t\beta_t\), and the
three negative sums are \(-\sum_t\mu_t^k g_t^k f_t^2\).
Using \(d_t^k=-\mu_t^k\), this is the quadratic form of
\eqref{eq:H2}. Polarization recovers the matrix action, which is
symmetric with respect to the sheet inner product.

\paragraph{Positivity and weighted kernel estimates}\label{par:inverse-kernel-estimates}
The estimates must retain the mass factors that will cancel the
denominators in \(P_k\). First, \(B^k\) is positive semidefinite. The
scalar continuation is jointly convex in its starting field and its
subsequent mass profile. Indeed, for a fixed bounded adapted control,
the controlled terminal field is affine in these variables, the
terminal value is convex, and the control cost is linear in the profile.
Taking the supremum preserves convexity; completing the square verifies
this control representation (PE, Section~5). The same calculation defines
the scalar PDE and control formula for every bounded nonnegative
profile, without imposing monotonicity. A strictly positive step profile
therefore admits both signs of every bounded profile direction, so its
full Hessian is nonnegative.
Approximate impulses by such directions, and only then average the
quadratic form over the prefix. The coefficients \(m_t\), \(t<k\),
in \(\beta_t\) are held fixed during this differentiation.

Here are explicit kernels for the impulse limit and its estimates. Put
\(u(y)=V_z(y,Z_y)\),
\(L^{\,y}(s)=\mathbf E_s[u(y)^2]/2\) for \(s\le y\), and let
\(L_z^{\,y}(s)\) denote its derivative with respect to the field at
time \(s\), evaluated at \(Z_s\). The subscript \(z\) is a field
derivative; the separate symbol \(z'\) below is a time. Order the two
evaluation times as \(y\le z'\), and keep the starting time
\(A=K_{k-1}\) visible. If both times precede \(A\), the kernel is
\(\mathbf E[u(y)u(z')\chi(A)]\). If
\(y\le A\le z'\), it is \(\mathbf E[u(y)L_z^{\,z'}(A)]\).
If both times follow \(A\), it is
\[
 \mathbf E[u(y)L_z^{\,z'}(y)]
       +\int_A^y \alpha(s)\mathbf E[L_z^{\,y}(s)L_z^{\,z'}(s)]\,ds .
\]
These formulas come from the second derivative of the diffusion PDE.
For bounded directions \(\eta,\eta'\), its mixed source is
\(\eta u D_zD_\alpha V[\eta']+
\eta' u D_zD_\alpha V[\eta]+
\alpha D_zD_\alpha V[\eta]D_zD_\alpha V[\eta']\).
Letting the directions concentrate at the two indicated times gives
the kernels above. When the two times coincide, the first term has
the common limit \(\mathbf E[u(y)^2\chi(y)]\).
The formulas also agree at \(y=A\) and \(z'=A\), since
\(L_z^{\,A}(A)=u(A)\chi(A)\). Thus no additional boundary term
appears when an evaluation time crosses \(A\).
For \(k=0\), the initial field is \(Z_0=0\). Field parity, bounded
derivatives, and the Brownian moment bound at the earlier time give
\(|B^0_{rt}|\lesssim C(1+K_{R+1})^C K_{\min(r,t)}\) when that time
is small.

The kernels \(B^k,g^k\) are bounded and Lipschitz in the evaluation
times, with polynomial constants. The spatial derivatives needed for
this claim are uniformly bounded: the source and attachment rules
express \(u,\chi\), their products, and the marked conditional
expectations as finite signed spin moments, even after any fixed number
of common-field derivatives.

For the kernel with both times after \(A\), varying the later time
uses \(\partial_{z'}L^{\,z'}(s)=\mathbf E_s\chi(z')^2/2\).
Varying the earlier time also differentiates \(u(y)\) and
\(L_z^{\,z'}(y)\); the latter has drift
\(-\alpha\chi L_z^{\,z'}\). For an evaluation time before \(A\),
condition the later factors on that time and apply It\^o's product rule
with \(u(y)\). The bounded derivatives just described control all
terms. The same argument applies to \(g^k\), using the drift
\(-\alpha\chi^2\) of \(\chi\).

Moving the starting time from \(K_{k-1}\) to \(K_k\) gives the
stronger estimate
\[
 \|B^{k+1}-B^k\|_\infty+\|g^{k+1}-g^k\|_\infty
    \lesssim_{\rm pol} p_k(K_k-K_{k-1}).
\]
On this interval \(\alpha=p_k\). In the three cases above, the
derivative with respect to \(A\) is respectively
\(-\alpha\mathbf E[u(y)u(z')\chi(A)^2]\),
\(-\alpha\mathbf E[u(y)\chi(A)L_z^{\,z'}(A)]\), and
\(-\alpha\mathbf E[L_z^{\,y}(A)L_z^{\,z'}(A)]\).
They remain bounded when a marker crosses \(A\), where the formulas
match. Integration over the interval supplies the factor
\(p_k(K_k-K_{k-1})\).

For the child kernels, set \(a_i=\mathbf E_t\chi_{i-1}\), and write
\(\chi(s)=V_{zz}(s,Z_s)\) for the continuous scalar-time process.
Its supermartingale drift gives
\[
 a_i-a_{i+1}
 =p_i\int_{K_{i-1}}^{K_i}\mathbf E_t\chi(s)^2\,ds.
\]
Thus a mixed difference of \(g_t^{il}=\mathbf E[a_i a_l]\) is
the expectation of the product of the two differences. This proves
the stated single and mixed bounds. The same product bound holds for
\((1-b_tg_t^{il})^{-1}\), as used above: interpolate the two
factors \(a_i,a_l\) linearly and differentiate
\(F(z)=(1-b_tz)^{-1}\) twice. The interpolated kernel is bounded
above by \(\mathbf E\chi_t^2\), so the gap bounds both derivatives
of \(F\). This includes \(i=l\); flat clock intervals follow by
continuity. Positivity and the same supermartingale drift give
\(0\le g_t^k,g_t^{il}\le\mathbf E\chi_t^2\), completing the
scalar estimates.
\end{proof}
